\section{Hamilton-cycle entropy with a fixed spectral gap}
\label{sec:enumeration}
This section proves a counting result under a fixed spectral gap;
no assumption that the nontrivial spectral ratio tends to zero
is made. For comparison, the bipartite enumeration theorem of
Byrne--Tait~\cite[Theorem 7]{BT}, extending Krivelevich~\cite{Krivelevich},
requires \(d/\lambda\geq(\log N)^{1+\epsilon}\).
Our explicit polylogarithmic degree threshold permits a
nontrivial ratio arbitrarily close to 1.

For a simple graph \(G\), let \(h(G)\) count unoriented Hamilton
cycles, with cyclic rotations identified.
Let \(J_m\) denote the \(m\)-by-\(m\) all-ones matrix.

\begin{theorem}[Full exponential Hamilton entropy at fixed gap]
\label{thm:hamiltonentropy}
Fix \(0<\delta<1\).
Let \(G=G_N\) be a sequence of simple balanced bipartite
graphs on \(N=2m\) vertices, with biadjacency matrix \(B\).
Suppose there are positive numbers \(D=D_N\) such that
\[
 \begin{gathered}
 D\geq(\log N)^{K_0},\qquad K_0=2\cdot10^7,\\
 \max_v|\deg_G(v)/D-1|=o(1),\qquad
 \left\|B/D-J_m/m\right\|\leq1-\delta .
 \end{gathered}
 \tag{43}
\]
Then
\[
 \log h(G)=N(\log D-1)+o(N),\qquad
 h(G)=(D/e)^N\exp(o(N)).                                    \tag{44}
\]
The logarithmic assertion includes eventual positivity.
The error is uniform along any sequence meeting the displayed
degree defect and fixed parameters.
\end{theorem}

We first supply the random-restriction estimate used in its proof.
This is a standard matrix-concentration mechanism, with
related formulations in~\cite[Lemmas 2.2--2.4]{Muyesser}
and the earlier work credited there.
We give the particular estimate needed below, without
claiming a new sampling principle or independence of
arithmetic entries.

\begin{lemma}[Two independent coordinate restrictions]
\label{lem:randomrestriction}
Let \(E\) be an \(m\)-by-\(m\) real or complex matrix, and put
\[
 v=\max\left\{\max_i\|E_{i,*}\|_2,\max_j\|E_{*,j}\|_2\right\},
 \qquad k=\max_{ij}|E_{ij}|,\qquad \ell=\log(4m).
\]
Choose independent row and column sets by Bernoulli sampling
with probability \(p\), with diagonal selectors \(R_X,R_Y\).
For every fixed \(A>0\), with probability at least \(1-m^{-A}\),
\[
 \|R_XER_Y\|
 \leq p\|E\|+C_A\bigl(v\sqrt{p\ell}+k\ell\bigr).             \tag{45}
\]
For independent uniform sets of common prescribed size \(b\),
take \(p=b/m\); the same bound fails with probability at most
\((m+1)^2m^{-A}\).
\end{lemma}
\begin{proof}
Write \(c_j\) for the columns of \(E\).
For the first column restriction the centered rank-one sum has
increments of norm at most \(v^2\) and matrix variance at most
\[
 p\sum_j\|c_j\|_2^2c_jc_j^*
             \preceq pv^2EE^*.
\]
Matrix Freedman applied to these independent increments gives,
with arbitrarily prescribed polynomially high probability,
\[
 \|ER_YE^*-pEE^*\|
 \leq C_A\bigl(v\|E\|\sqrt{p\ell}+v^2\ell\bigr).
\]
Taking square roots, and increasing \(C_A\), yields
\[
 \|ER_Y\|\leq\sqrt p\|E\|+C_Av\sqrt\ell.                     \tag{46}
\]
For each row, the sum
\(\sum_j\one_{j\in Y}|E_{ij}|^2\) has mean at most \(pv^2\),
bounded summands \(k^2\), and scalar variance at most
\(pk^2v^2\).
The scalar version of the same inequality and a union bound
over rows give
\[
 \max_i\|(ER_Y)_{i,*}\|_2\leq\sqrt p\,v+C_Ak\sqrt\ell.        \tag{47}
\]
Condition on \(Y\), and apply the rank-one matrix argument to
the independent row restriction of \(Z=ER_Y\).
It gives
\[
 \|R_XZ\|\leq\sqrt p\|Z\|
                   +C_A\sqrt\ell\max_i\|Z_{i,*}\|_2.
\]
Substitute (46)--(47), retaining both error terms, to obtain
(45). Choosing the failure exponents a little larger before
the three union bounds gives the stated exponent \(A\).

If \(p=b/m\), the event that a Bernoulli set has size exactly
\(b\) has probability at least \(1/(m+1)\): \(b\) is a mode
of its binomial distribution.
Conditional on that event the set is uniform.
The two independent conditioning events therefore cost at
most \((m+1)^2\) in failure probability.
The cases \(b=0,m\) follow directly, or can be omitted in the
application where \(0<b<m\).
\end{proof}

\begin{lemma}[Regular expanding factors in random blocks]
\label{lem:randomblockfactor}
Assume the uniform-degree and centered-norm hypotheses
of (43), without imposing its global lower bound on \(D\).
Let \(X,Y\) be independent uniform subsets of the two parts,
each of size \(s\), and put \(p=s/m\).
If \(pD/\log N\to\infty\), then, with probability
\(1-N^{-A}\) for every fixed prescribed \(A\), the block
\(G[X,Y]\) has a regular spanning factor of degree
\[
 d_s=(1-o(1))pD
\]
and a fixed positive nontrivial gap depending only on \(\delta\).
All vertex degrees before factoring are \((1+o(1))pD\).
If the host is exactly \(D\)-regular, the factor can be chosen
so that its degree differs from \(pD\), and its discarded
maximum degree is bounded, by
\[
 C_{\delta,A}\bigl(\sqrt{pD\log N}+\log N\bigr).
 \tag{48}
\]
The assertions hold simultaneously for polynomially many
specified block pairs of a random ordered partition.
A common factor degree may be used for pairs with the same size.
\end{lemma}
\begin{proof}
Use \(E=B/D-J_m/m\).
The degree condition in (43) and \(D\leq(1+o(1))m\)
give
\[
 v=O(D^{-1/2}),\qquad k=O(D^{-1}).
\]
Lemma~\ref{lem:randomrestriction}, conditioned on the two
prescribed sizes, gives
\[
 \left\|\frac{B[X,Y]}{pD}-\frac{J_s}{s}\right\|
 \leq1-\delta+
        O_A\left(\sqrt{\frac{\log N}{pD}}
                        +\frac{\log N}{pD}\right).
 \tag{49}
\]
Sampling neighbours without replacement gives, uniformly for
all vertices in the original graph,
\[
 \deg(v,Y)\ \hbox{or}\ \deg(v,X)
 =pD\left(1+O\left(\varepsilon_N+
       \sqrt{\frac{\log N}{pD}}+\frac{\log N}{pD}\right)\right),
 \tag{50}
\]
where \(\varepsilon_N=\max_v|\deg(v)/D-1|\).
The usual scalar concentration inequality follows also by
Bernoulli sampling and the same conditioning argument.
The constants can be chosen for any polynomial failure exponent.

For large \(N\), (49) is at most \(1-\delta/2\) and the
degree defect in (50) tends to zero.
The Laplacian argument proving (28), with the bipartite
constant block \(J_s/s\), now gives a fixed positive cut
constant, for example \(g=\delta/8\).
Lemma~\ref{lem:regularfactor} extracts a factor of degree
\((1-o(1))pD\), deleting \(o(pD)\) edges per vertex.
The discarded adjacency has the same operator bound;
(49) transfers a fixed positive nontrivial gap to the factor.
All errors depend on the displayed deterministic defect and
on \(pD\), so one may take their common upper bound before
choosing the integer factor degree.
For an exactly regular host, (50) has no \(\varepsilon_N\)
term. The explicit loss in (22)--(23), including its floor,
is then at most a constant times
\(\sqrt{pD\log N}+\log N\), which proves (48).
The marginals of any two opposite-part blocks in a random
ordered partition are independent uniform subsets.
A union bound with a larger prescribed exponent proves the
simultaneous assertion; independence between different block
pairs is not required.
\end{proof}

\begin{proof}[Proof of Theorem~\ref{thm:hamiltonentropy}]
We first regularize the host before sampling.
The centered norm in (43) and its vanishing degree defect
give a fixed positive cut constant by the Laplacian argument
used in (28).
Lemma~\ref{lem:regularfactor} gives an exactly \(d\)-regular
spanning subgraph \(H\), where
\[
 d=(1-o(1))D,\qquad
 \|A_G-A_H\|=o(D).
\]
Consequently its biadjacency satisfies
\(\|B_H/d-J_m/m\|\leq1-\delta/2\) for all large \(N\).
Also \(d\geq\tfrac12(\log N)^{K_0}\) eventually.
For the lower bound work in this exact regular host;
every resulting cycle belongs to the original graph.
The difference \(\log d-\log D=o(1)\) will cost only \(o(N)\).

Take
\[
 t=\lfloor d^{1/4}\rfloor,\qquad
 b=\left\lfloor\frac{m-m/(\log N)^3}{t}\right\rfloor,\qquad
 r=m-tb.
\]
These parameters satisfy
\[
 t\asymp d^{1/4},\quad b\asymp N d^{-1/4},\quad
 r=\frac{m}{(\log N)^3}+O(d^{1/4}),\quad
 (r+b)\log N=o(N),\quad t\log N=o(N).                        \tag{51}
\]
Partition the first part into ordered sets
\(A_1,\ldots,A_t,R_A\), of sizes \(b,\ldots,b,r\),
and independently the other part into
\(B_1,\ldots,B_t,R_B\) with the same sizes.
The number of ordered partition pairs is
\[
 {\cal P}=\left(\frac{m!}{r!(b!)^t}\right)^2.
\]
Stirling's formula and (51) give, with \(p=b/m\),
\[
 \log{\cal P}=2tb\log(1/p)+o(N).                            \tag{52}
\]

Require the regular-factor conclusions in every adjacent
pair of the alternating list
\[
 A_1,B_1,A_2,B_2,\ldots,A_t,B_t.
\]
Their sampled degree is \(pd\asymp d^{3/4}\),
so it exceeds \((\log N)^{3K_0/4}\) up to a fixed factor.
Require the same conclusions on
\[
 K=H[R_A\cup A_1,\ R_B\cup B_t],\qquad
 n_K=2(r+b),\qquad p_K=(r+b)/m .
\]
The augmented reservoir has independent uniform part sets
of size \(r+b\), and \(p_Kd\asymp d/(\log N)^3\).
In addition require the scalar endpoint concentration
\[
 |\deg_H(v,A_1)-pd|,\ |\deg_H(v,B_t)-pd|
       \leq C\bigl(\sqrt{pd\log N}+\log N\bigr)              \tag{53}
\]
for every vertex \(v\) of the appropriate opposite part.
This holds by sampling without replacement in the exact
regular host.
There are only \(2t\) block requirements and \(2N\) scalar
requirements. Lemma~\ref{lem:randomblockfactor} and scalar
concentration, with suitably larger failure exponents,
therefore hold for a proportion \(1-o(1)\) of all partitions.
Their events depend on the partition, not on any matching.

Fix a good partition.
For every consecutive block choose a factor with the same
degree \(d_0=(1-o(1))pd\).
The van der Waerden permanent bound gives at least
\[
 d_0^b\,\frac{b!}{b^b}\geq(d_0/e)^b
\]
perfect matchings in each.
This established permanent inequality, together with
Bregman's upper bound used below, is recorded in
\cite[Section 2.4]{BT}.
Choose these \(2t-1\) matchings as independent combinatorial
choices. Their union covers the \(2tb\) nonreservoir vertices
by \(b\) paths, each with \(2t\) vertices.
The first endpoints are \(A_1\), the last endpoints \(B_t\).
The layer structure prevents cycles, and each resulting path
cover determines all its chosen matchings.

We verify the source hypotheses for closing every cover.
Choose in \(K\) an exact regular factor \(K'\) of degree \(d_K\)
using (48), so that
\[
 |d_K-p_Kd|,\quad
 \max_v\deg_{K-K'}(v)
       \leq C_\delta\bigl(\sqrt{p_Kd\log N}+\log N\bigr).
 \tag{54}
\]
It has fixed positive nontrivial gap.
The endpoint sets can be too large for the arbitrary-set
bound in Lemma~\ref{lem:smallendpoints}; their distribution
must be retained instead.
Combining (53)--(54), for every \(v\) of the opposite part,
\[
 \left|\deg_{K'}(v,A_1)-\frac b{r+b}d_K\right|,
 \quad
 \left|\deg_{K'}(v,B_t)-\frac b{r+b}d_K\right|
 \leq C_\delta\bigl(\sqrt{p_Kd\log N}+\log N\bigr).           \tag{55}
\]
Indeed the target before factoring is \(pd=(b/(r+b))p_Kd\).
The three errors are the original endpoint fluctuation,
the actual number of edges lost at the vertex, and
\((b/(r+b))|d_K-p_Kd|\).
Each is bounded by the right side, since \(p\leq p_K\).
Thus no endpoint compensation is erased by regularization.

The two distinct endpoint sets, both of size \(b\), satisfy
the set and multiset bi-uniformity condition (41), now on
\(K'\), with
\[
 \alpha=C_\delta\left(
       \sqrt{\frac{\log N}{p_Kd}}+\frac{\log N}{p_Kd}\right).
\]
Take
\[
 c=1/5,\qquad \gamma=(\log n_K)^{-3},\qquad d'=0 .
\]
The degree satisfies
\(d_K\asymp d/(\log N)^3\gg(\log N)^{K_0-3}\).
In particular
\[
 \alpha d_K^c
 \ll_\delta(\log N)^{1/2}d_K^{c-1/2}
                       +(\log N)d_K^{c-1}\longrightarrow0.
\]
Thus \(\alpha\leq d_K^{-c}\) for all large \(N\).
For the number \(b\) of prescribed pairs, charge the full range:
\[
 \frac{b}{d_K^{-c/900}n_K}
 \ll (\log N)^3 d^{-1/4+c/900}\longrightarrow0.
 \tag{56}
\]
The limit follows from \(K_0(1/4-1/4500)>3\).
The fixed positive spectral gap gives a fixed positive
edge-expansion constant, eventually larger than \(\gamma\).
Also
\[
 \gamma^{-10^6/c}=(\log n_K)^{15\cdot10^6}
                   =o(d_K),
\]
since \(K_0-3>15\cdot10^6\).
The remaining source conditions are
\(0<c<1/4\), \(\gamma<(\log n_K)^{-2}\), equal endpoint
cardinalities, \(d'=0\leq d_K^{1-c}\), and large \(n_K\)
relative to fixed \(c\); all hold.

Label the long paths \(Q_1,\ldots,Q_b\), with endpoints
\(a_i\in A_1,z_i\in B_t\), and apply (42) to
\((a_{i+1},z_i)\), taking indices cyclically.
It supplies disjoint connecting paths with interiors
exactly \(R_A\cup R_B\).
Together with the long paths, in the prescribed cyclic order,
these form one Hamilton cycle.
For all large \(N\), \(t\geq2\), so no long-path edge has both
ends in \(K\); the two edge collections do not overlap.

Choose one such Hamilton cycle for each configuration by a
fixed deterministic rule.
For a fixed output cycle \(C\), its long-path cover in any
preimage is a subset of its \(N\) edges with \(2tb-b\) edges.
The number of possibilities is at most
\[
 \binom N{\,2r+b\,}=\exp(o(N)).                              \tag{57}
\]
Each subset determines its entire configuration:
its isolated vertices are the reservoir, and each nontrivial
component is one long path.
Orient each path from its endpoint in the first part;
successive positions fix all ordered layers and matchings.
No permutation of paths or extra orientation factor is
left uncharged.

The number of good configurations is at least
\[
 (1-o(1)){\cal P}\,(d_0/e)^{b(2t-1)}.
\]
Its logarithm, by (51)--(52), is
\[
 2tb\log(1/p)+(2tb-b)(\log(pd)-1)+o(N)
                 =N(\log d-1)+o(N).
\]
The missing \(2r+b\) matching edges cost
\(O((r+b)\log N)=o(N)\).
The ordered-partition entropy cancels the entire layer
degree loss. Division by (57), followed by
\(\log d=\log D+o(1)\), proves the required lower bound
for the original graph.

For its upper bound, the full adjacency permanent counts
all directed cycle covers and includes both orientations
of every Hamilton cycle.
Bregman's inequality gives
\[
 2h(G)\leq\operatorname{per}(A_G)
            \leq\prod_v(\deg_G(v)!)^{1/\deg_G(v)}.
\]
Stirling and (43) imply
\[
 \log\operatorname{per}(A_G)
 \leq\sum_v(\log\deg_G(v)-1)
       +O\left(\frac{N\log D}{D}\right)
 =N(\log D-1)+o(N).
\]
This completes the proof with its explicit polylogarithmic
degree threshold.
\end{proof}

\begin{corollary}[The exponential asymptotic for prime circles]
\label{cor:primecircleentropy}
Fix an even \(q\) and a nonempty allowed unit set \(R\).
If \(g_R=2\), then along the even integers
\[
 h(G_N(q,R))
 =\left(\frac{\sigma_R N}{e\log N}\right)^N\exp(o(N)).
 \tag{58}
\]
If \(g_R>2\), the count is zero for all \(N\geq g_R\).
In particular the ordinary prime-sum graph satisfies
\[
 \log h(G_N)
   =N\log N-N\log\log N-N+o(N).
 \tag{59}
\]
\end{corollary}
\begin{proof}
The norm estimate (40) and the stochastic cross-block
spectral gap imply
\[
 \left\|B_N^R/D-J_m/m\right\|
                    \leq1-\delta_R+o(1),
 \qquad D=\sigma_R N/\log N .
\]
The uniform degrees are (40), and
\(D\geq(\log N)^{K_0}\) once \(N\) is sufficiently large.
Apply Theorem~\ref{thm:hamiltonentropy} with any smaller
fixed positive gap.
The disconnected case is the exact obstruction in
Theorem~\ref{thm:filterclassification}.
For \(q=2,R=\{1\}\), the relative prime density is one,
giving (59).
\end{proof}

\begin{remark}[Precision of the enumeration]
The factor \(\exp(o(N))\) is part of the assertion.
This determines the exponential base, including its constant
\(e^{-1}\); it is not a relative-error \(1+o(1)\) formula for
the total count.
It leaves, for example, any bounded singular-series factor
and finer subexponential terms undetermined.
\end{remark}

\begin{corollary}[An entropy and concentration inverse for kernel supports]
\label{cor:entropyinverse}
Fix \(0<\delta<1\), and let \(K_0=2\cdot10^7\).
There is \(a_\delta>0\) such that, for sufficiently large even
\(N\), a stochastic bipartite kernel \(P\) with nontrivial
gap at least \(\delta\) and
\[
 \alpha=\max_{ij}P_{ij}\leq a_\delta L^{-(K_0+2)}             \tag{60}
\]
has Hamiltonian support.
More quantitatively, with \(c_\delta\) from the finite
rounding theorem and
\[
 d=\left\lfloor\frac{c_\delta}{2\alpha L^2}\right\rfloor,
 \qquad
 \log h(\supp P)\geq N(\log d-1)-o(N),                       \tag{61}
\]
the displayed lower bound holds along every sequence
satisfying (60), with fixed \(\delta\).
Consequently a non-Hamiltonian support with this fixed gap
forces \(\alpha>a_\delta L^{-(K_0+2)}\).
\end{corollary}
\begin{proof}
Stochasticity on equal parts gives \(\alpha\geq1/m=2/N\).
Decrease \(c_\delta\) if necessary so that \(c_\delta\leq1\),
and choose \(a_\delta\leq c_\delta/8\).
Then the degree in (61) is at most \(m\) and, for all
sufficiently large \(N\), is at least
\(2L^{K_0}\), hence at least \((\log N)^{K_0}\).
It obeys \(d\alpha L^2\leq c_\delta/2\).
The full rounding error in (1) is bounded by
\[
 C_\delta\left(\frac{c_\delta}{2L}
       +\sqrt{\frac Ld}+\frac Ld\right)=o(1).
\]
The resulting simple regular graph \(H\) therefore has
nontrivial gap at least \(\delta/2\), and its cross block
satisfies (43) with degree \(d\).
Theorem~\ref{thm:hamiltonentropy} gives
\(\log h(H)=N(\log d-1)+o(N)\).
Since \(H\) is contained in the support graph, (61) follows.
The inverse is the contrapositive of its positivity conclusion.
The \(o(N)\) term denotes a uniform vanishing exponential
loss under the displayed hypotheses, not a numerical
starting threshold.
\end{proof}
